% Part: normal-modal-logic
% Chapter: filtrations
% Section: euclidean-filtrations

\documentclass[../../../include/open-logic-section]{subfiles}

\begin{document}

\olfileid{nml}{fil}{euc}

\olsection{ترشيحات النماذج الإقليدية}

أما النماذج الإقليدية، وحدها أو مع التسلسل، فلا تكفي فيها طريقة
\olref[acc]{sec}. وبيان ذلك النموذج المرسوم في~\olref{fig:ser-eucl}،
فهو متسلسل وإقليدي. خذ~$\Gamma=\{p,\Box p\}$؛ فإذا رُشّح النموذج
عبر~$\Gamma$ اتحدت الفئتان~$[w_1]=[w_3]$؛ فإن~$w_1$ و$w_3$
هما وحدهما، من بين أزواج العوالم المختلفة، يتفقان على~$\Gamma$.
ويجب كذلك أن تنتقل أسهم~$\mModel{M}$ إلى كل ترشيح، على نحو ما في
\olref{fig:ser-eucl2}. لكن النموذج الناتج ليس إقليديًا، ولا يمكن جعله
كذلك بزيادة أسهم مع إبقائه ترشيحًا. فإن الإقليدية توجب سهمين متعاكسين
بين~$[w_2]$ و$[w_4]$، ثم بين~$[w_2]$ و$[w_5]$.
وهذا ممتنع، لأن~$\Box p$ صادقة عند~$w_2$، و$p$ كاذبة عند~$w_5$.

\begin{figure}[htpb]
  \centering
  \begin{tikzpicture}[modal]
    \node[world] (w1)
          [label={left: \mFalse{p}},
            label={below: $\mSat{{}}{\Box p}$}] {$w_1$};
    \node[world] (w2)
          [label={right: \mTrue{p}},
            label = {below: $\mSat{{}}{\Box p}$},
            right=of w1] {$w_2$};
    \draw[->] (w1) to (w2);
    \draw[reflexive above] (w2) to (w2);
    \node[world] (w3)
          [label={left: \mFalse{p}},
            label={below: $\mSat{{}}{\Box p}$},
            below=of w1] {$w_3$};
    \node[world] (w4)
          [label={right:\mTrue{p}},
            label={below: $\mSat/{{}}{\Box p}$},
            right=of w3] {$w_4$};
    \draw[->] (w3) to (w4);
    \draw[reflexive above] (w4) to (w4);
    \node[world] (w5)
          [label={right: \mFalse{p}},
            label=below:$\mSat/{{}}{\Box p}$,
            right=of w4] {$w_5$};
    \draw[->, bend left] (w4) to (w5);
    \draw[->, bend left] (w5) to (w4);
    \draw[reflexive above] (w5) to (w5);
  \end{tikzpicture}
  \caption{نموذج متسلسل وإقليدي.}
  \ollabel{fig:ser-eucl}
\end{figure}

\begin{figure}[ht]
  \centering
  \begin{tikzpicture}[modal,every text node part/.style={align=center}]
    \node[world] (w1)
          [label={left: \mFalse{p}},
            label={right:$[w_1]=[w_3]$},
            label={below: $\mSat{{}}{\Box p}$}] {$[w_1]$};
    \node[world] (w2)
         [label={right: \mTrue{p}},
           label = {below: $\mSat{{}}{\Box p}$},
           right=of w1, yshift=1.5cm] {$[w_2]$};
    \draw[->] (w1) to (w2);
    \draw[reflexive above] (w2) to (w2);
    \node[world] (w4)
         [label={right:\mTrue{p}},
           label={below: $\mSat/{{}}{\Box p}$},
           right=of w1,yshift=-1.5cm] {$[w_4]$};
    \draw[->] (w1) to (w4);
    \draw[reflexive above] (w4) to (w4);
    \node[world] (w5)
         [label={right: \mFalse{p}},
             label=below:$\mSat/{{}}{\Box p}$,
             right=of w4] {$[w_5]$};
    \draw[->, bend left] (w4) to (w5);
    \draw[->, bend left] (w5) to (w4);
    \draw[reflexive above] (w5) to (w5);
  \end{tikzpicture}
  \caption{ترشيح النموذج في~\olref{fig:ser-eucl}.}
  \ollabel{fig:ser-eucl2}
\end{figure}

فلا يكفي، للحصول على ترشيح إقليدي بهذه الطريقة، إغلاق مجموعة اعتباطية
$\Gamma$ تحت أخذ ال!!{formula}s الجزئية. وإنما ننظر في مجموعات
$\Gamma$ \emph{مغلقة جهويًا}، بالمعنى المبين في
\olref[pre]{defn:modallyclosed}. غير أن هذه المجموعات، متى كانت غير
خالية، لا نهائية؛ فلا يُستفاد منها مباشرةً إثبات خاصية النموذج المنتهي،
ولا قابلية النظام المقابل للقرار.

\begin{thm}\ollabel{thm:modal-closed-filt}
  لتكن~$\Gamma$ مغلقة جهويًا، وليكن~$\mModel{M}=\tuple{W,R,V}$، وليكن
  $\mModel{M^*}=\tuple{W^*,R^*,V^*}$ الترشيح الأخشن لـ~$\mModel{M}$
  عبر~$\Gamma$.
  \begin{enumerate}
  \item إذا كان~$\mModel{M}$ متعديًا، كان~$\mModel{M^*}$ متعديًا.
  \item إذا كان~$\mModel{M}$ إقليديًا، كان~$\mModel{M^*}$ إقليديًا.
  \item إذا كان~$\mModel{M}$ متناظرًا، كان~$\mModel{M^*}$ متناظرًا.
  \end{enumerate}
\end{thm}

\begin{proof}
\begin{enumerate}
  \item لما كان~$\mModel{M^*}$ الترشيح الأخشن، كانت~$R^*[u][v]$
    مكافئة لـ~$C_1(u,v)$ بالتعريف. فالتعدي يثبت بأن نفرض~$C_1(u,v)$
    و$C_1(v,w)$، ثم نستنتج~$C_1(u,w)$.
    \iftag{prvBox}{لتكن~$\Box !A \in \Gamma$، وافترض
      $\mSat{M}{\Box !A}[u]$؛ عندئذ
      $\mSat{M}{\Box\Box!A}[u]$، لأن~\Ax{4} صالحة في جميع النماذج
      المتعدية؛ ولأن~$\Box\Box!A \in \Gamma$ بالإغلاق، ينتج من
      $C_1(u,v)$ أن~$\mSat{M}{\Box!A}[v]$، ثم ينتج من~$C_1(v,w)$ أن
      $\mSat{M}{!A}[w]$. }{}\iftag{prvDiamond}{لتكن
      $\Diamond !A \in \Gamma$، وافترض $\mSat{M}{!A}[w]$؛
      عندئذ~$\mSat{M}{\Diamond !A}[v]$ بحسب
      $C_1(v,w)$، لأن~$\Diamond !A \in \Gamma$ بالإغلاق الجهوي. ومن
      $C_1(u,v)$ نحصل على~$\mSat{M}{\Diamond\Diamond !A}[u]$، لأن
      $\Diamond\Diamond!A \in \Gamma$ بالإغلاق الجهوي. ولأن
      \Ax{4_\Diamond} صالحة في جميع النماذج المتعدية، يكون
      $\mSat{M}{\Diamond!A}[u]$.}{}
    % تصحيح برهاني (0459-I1): صُرّح بفرضي العضوية اللازمين لإثبات الشرط C_1.
  \item تمرين. استعمل كون كل من~\Ax{5} و$\Ax{5_\Diamond}$ صالحة في
    جميع النماذج الإقليدية.
  \item تمرين. استعمل كون~\Ax{B} و$\Ax{B_\Diamond}$ صالحة في جميع
    النماذج المتناظرة.
\end{enumerate}
\end{proof}

\begin{prob}
  أكمل برهان~\olref[nml][fil][euc]{thm:modal-closed-filt}.
\end{prob}

\end{document}
